% Français 9115, batch 11 — Gallica views 201–220.
% Pass: Opus 5 (claude-opus-5), 2026-09-19 — first pass, unchecked against
% the views by a human.
%
% What the twenty views hold. The batch turns at view 206. The Latin fair
% copy of Euler's « Investigatio motuum quibus laminæ et virgæ elasticæ
% contremiscunt » (Acta Petropolitana pro anno 1779, Pars prior), which has
% run since view 144, ends at view 206 — the copy's page (161), closed by
% the French word « fin. » and a paraph. Everything after it is Sophie
% Germain's own French: a title leaf, then two drafts of a « Remarque(s)
% sur le mémoire d'Euler », the second of which begins over again with its
% own title and its own leaf numbering. Twelve views carry text.
% View 204 is not transcribed because it photographs page (159) a second
% time (view 203) — same text, same ink number 99, only the framing differs.
% Seven views are blank versos carrying nothing but the show-through of
% their own recto: 207, 209, 211, 213, 215, 217, 219.
%
%   v. 201–206  End of the Euler copy, pages (157) to (161). § 50 finishes
%               the coefficients for the rod pinned at its middle
%               ($\alpha'=0$, $\beta'=0$, $\gamma'=-\gamma$, $\delta'=0$),
%               then § 51 takes the second factor, tang $\frac12\omega$ =
%               $(e^\omega-1)/(e^\omega+1)$, the same equation as case V,
%               giving $\omega = 5\pi/2, 9\pi/2 \dots$ and sounds two octaves
%               above case V. Scholion § 52: the rod naturally curved, the
%               axis taken along the curve $EXF$ itself, the same differential
%               equation and the same integral, « Tab III fig 9 ». § 53–55:
%               the closing case, a rod whose natural figure closes on itself
%               (« Tab. III fig 10 ») — de sonis annulorum elasticorum: all
%               four coefficients but $\gamma$ vanish, $\omega = i\pi$, the
%               sounds go as the squares, and such rings, like discs and
%               bell-shaped basins, sound purer than strings. The copy ends
%               there, « fin. »
%   v. 208      A title leaf in her hand and nothing else:
%               « Remarques sur la lame elastique ».
%   v. 210–216  « Remarque sur le mémoire d'Euler … » — a first draft, her
%               leaves numbered 2, 3, 4 at the top left (the first leaf's
%               figure is a blot). She restates § 47–51 in French, then
%               objects: « Il me semble étonnant qu'Euler ait pû croire que
%               cette seconde solution satisfait a laquestion », since his
%               hypothesis is general, the condition of communication between
%               the parts never enters the calculation, and for every
%               $\lambda$ but $\frac12$ the second factor leads only to
%               $\omega = 0$, that is, to rest. She then sets out to prove
%               that for any $\lambda$ that is an exact fraction of unity the
%               final equation always admits a solution belonging to case IV,
%               and transforms it in six steps into a form homogeneous in
%               $\lambda$ and $1-\lambda$.
%   v. 218–220  « Remarques sur le memoire d'Euler … » — a second, fuller
%               draft, beginning again with the title and with its own leaf
%               numbering 1, 2 at the top middle. It opens on the history
%               (Euler had treated the question in the Methodus inveniendi
%               curvas maximi minimique gaudentes, 1744; nothing essential
%               has changed), announces that she will report this part of
%               the memoir « en y joignant les problèmes relatifs aux cas
%               qu'il a négligé de traiter », and then quotes Euler at
%               length in French, a guillemet at the head of every line,
%               through the Probleme and Solution of § 47. It runs on past
%               view 220.
%
% The numbers on the leaves, and why no \folio{} is written. Three kinds,
% all in ink. (1) The Latin copy's parenthesised page numbers at the head of
% the page, in the middle: (157) at view 201, (158) at 202, (159) at 203
% (and again on its second photograph, view 204), (160) at 205, (161) at
% 206. As in batches 8 to 10 these are very probably the pages of the
% printed volume being copied — an inference. (2) The ink series of the
% earlier batches at the top right of the page: 98 at view 201, 99 at 203
% (and 204), 100 at 206 — the pages with odd copy numbers, as before — then,
% on the French leaves, 101 at view 208, 102 at 210, 103 at 212, 104 at 214,
% 105 at 216, 106 at 218, 107 at 220, one per written recto without a gap.
% (3) Her own numbering of each draft's leaves: 2 at view 212, 3 at 214, 4
% at 216, at the top left (view 210, which should carry 1, has an inked
% triangular blot there that is not read); and, for the second draft, 1 at
% view 218 and 2 at view 220, at the top middle, each underlined. No
% pencilled foliation of the Bibliothèque was read anywhere in the batch, so
% no \folio{} is written, as in batches 1 to 10. Every number above was read
% on its view; none is inferred.
%
% Notation. In the Latin copy the letters are kept as in batches 8 to 10:
% $\alpha$, $\beta$, $\gamma$, $\delta$ (and primes) for the constants,
% $\omega = a/f$, $u = s/a$, $\zeta$ the phase, $C$ the arbitrary
% coefficient, $\lambda$ for $EL/EF$. The radical is written without a bar
% and is set \surd; at view 203 it carries a small index read as 4 and is
% set $\sqrt[4]{bcck}$, which draws a bar the leaf does not have. « sin. »,
% « cos. », « tang. » and « cot. » are written with and without their stop
% in the same line, and the pass has not tried to record every stop. In the
% French leaves she uses the same letters, and the ordinary signs of her
% time; « N\textsuperscript{o} » and « N. » both occur for Euler's section
% numbers and both are kept. Her orthography stands as written, including
% the words she runs together — « dela », « detous », « demanière »,
% « parconséquent », « aurepos », « cidessus », « delamême », « lemême »,
% « laquestion », « enquestion », « desorte », « aumilieu », « acause »,
% « apresent », « alavérité » — and her spellings « mouvemens »,
% « differens », « coëfficiens », « exclutif », « appuiées », « fesant »,
% « interromp », « pendul », « aprofondi », « naitre », « donnees ». Long
% horizontal strokes through a word are, here as in the Euler copy, the
% crossbar of its $t$ thrown back and not a deletion; only a separate,
% heavier stroke is written \struck{}. Where the calculation and the leaf
% disagree the leaf stands, with a note: view 202 (« sumto », « naturam »),
% view 214 (the marginal $\text{sin}\,\omega$), view 220 (« ce point ce
% point »).
%
% What could not be read. Two places are marked \ill{}: the word struck
% before « percipietur » at view 206, under a heavy erasure; the two-letter
% abbreviation struck at the head of the text of view 218. Readings offered
% with doubt are marked \uncertain{}: « hance » at view 202; the short word
% between « puisse » and « etre appuiée » at view 212, read « pu », which
% makes no sense; « 2d » before « N\textsuperscript{o} 35 » at view 210; the
% blotted sign before the last marginal fraction at view 214, read « 2 »;
% the page cited after « l'auteur » at view 218, read « p. 151 », whose
% middle figure is not settled between 3 and 5. Nothing was guessed to fill
% a gap.
%
% Nothing here was checked against any published transcription or edition
% of Euler's memoir or of Sophie Germain's papers.
\documentclass[11pt,a4paper]{article}
% The preamble shared by every transcription of the mathematicians' manuscripts in Gallica.
%
% It rests on one idea: **uncertainty compounds, it does not resolve.** A
% transcription that smooths over a doubtful word has destroyed the only thing
% it was made for — a reader must be able to tell, at every point, what was on
% the page and what was supplied. The apparatus macros below are therefore the
% critical apparatus, not decoration.
%
% The same commands are recognised by `scripts/render.mjs`, which produces the
% site's reading view, and by `scripts/tei.mjs`. A macro added here must be
% added there too, or rendering fails loudly — which is the wanted behaviour.
%
% Comments here are English, like the rest of the repository. The documents
% this preamble sets are French, and that is deliberate: see the skills.

\ProvidesPackage{germain}[2026/09/15 Transcriptions of mathematicians' manuscripts in Gallica]

\usepackage{amsmath,amssymb,amsthm}
\usepackage{mathrsfs}
% Kept from the parent preamble so the renderer's diagram support stays
% available; these manuscripts draw no commutative diagrams, and nothing here uses it.
\usepackage{tikz-cd}
\usepackage[english,french]{babel}
\usepackage{fontspec}

% --- Greek ------------------------------------------------------------------
% The text face stays fontspec's default, Latin Modern, which has no Greek:
% XeTeX drops every glyph it lacks with a line in the log and nothing on the
% page, and the Greek words the manuscripts quote vanished from the PDFs. So
% the Greek and Greek Extended blocks get a XeTeX character class of their own,
% and entering or leaving a run of it switches the family to CMU Serif — the
% same Computer Modern design, with polytonic Greek — and back. Only the
% family changes: a Greek word in \textit or \textbf keeps its shape and series.
% The transitions to and from every other class carry the tokens that class
% has towards an ordinary letter, so French spacing before « ; » or inside
% « » still holds after a Greek word. No grouping, so no brace or \endgroup
% can come out unbalanced; maths is left alone, since XeTeX inserts nothing
% there. Kept identical in `exercices.sty`.
\makeatletter
\newfontfamily\germainGreekfont{cmun}[
  Extension=.otf, NFSSFamily=germaingreek,
  UprightFont=*rm, ItalicFont=*ti, SlantedFont=*sl,
  BoldFont=*bx, BoldItalicFont=*bi, BoldSlantedFont=*bl]
\newXeTeXintercharclass\germain@greek
\def\germain@greekrange#1#2{%
  \@tempcnta=#1\relax
  \loop
    \XeTeXcharclass\@tempcnta=\germain@greek
    \ifnum\@tempcnta<#2\advance\@tempcnta\@ne\repeat}
\germain@greekrange{"0370}{"03FF}
\germain@greekrange{"1F00}{"1FFF}
\def\germain@greekprev{\rmdefault}
\def\germain@greekin{%
  \edef\germain@greekprev{\f@family}\fontfamily{germaingreek}\selectfont}
\def\germain@greekout{\fontfamily{\germain@greekprev}\selectfont}
\def\germain@greekjoin{%
  \XeTeXinterchartokenstate=\@ne
  \@tempcnta=\z@
  \loop
    \ifnum\@tempcnta=\germain@greek\else
      \XeTeXinterchartoks\germain@greek\@tempcnta=\expandafter{%
        \expandafter\germain@greekout\the\XeTeXinterchartoks\z@\@tempcnta}%
      \XeTeXinterchartoks\@tempcnta\germain@greek=\expandafter{%
        \the\XeTeXinterchartoks\@tempcnta\z@\germain@greekin}%
    \fi
    \ifnum\@tempcnta<4095\advance\@tempcnta\@ne\repeat}
% babel-french sets its own classes, and saves and restores the interchar
% state, each time a language is selected: join again after it does.
\addto\extrasfrench{\germain@greekjoin}
\addto\extrasenglish{\germain@greekjoin}
\AtBeginDocument{\germain@greekjoin}
\makeatother

\usepackage{geometry}
\geometry{a4paper,margin=25mm,marginparwidth=28mm}
\usepackage{marginnote}
\usepackage{xcolor}
\usepackage[normalem]{ulem}
\setlength{\emergencystretch}{3em}
\usepackage{hyperref}
\hypersetup{colorlinks=true,linkcolor=black,urlcolor=black,pdfborder={0 0 0}}

% --- Batch metadata ---------------------------------------------------------
% Taken as they stand by the rendering script, which shows them at the head of
% the reading view. They fix what the file claims to cover. The macro names are
% the parent project's, so that the scripts stay shared; \folder here means the
% volume's slug — `fr-9115` — and \pages counts Gallica views.

\newcommand{\folder}[1]{\gdef\grothFolder{#1}}
\newcommand{\batch}[1]{\gdef\grothBatch{#1}}
\newcommand{\pages}[2]{\gdef\grothFirst{#1}\gdef\grothLast{#2}}
\newcommand{\foldertitle}[1]{\gdef\grothFoldertitle{#1}}

% The holder's own shelfmark, printed as they write it: « Français 9115 ».
\newcommand{\shelfmark}[1]{\gdef\grothShelfmark{#1}}

% The dating, copied verbatim from the holder's catalogue — never inferred
% here. « XIXe siècle » is the BnF's; a pass that dates a leaf from its content
% says so in a \note{}, as its own inference.
\newcommand{\dating}[1]{\gdef\grothDating{#1}}

% The status notice, printed in the title block and repeated as a diagonal
% watermark on every page of the PDF. The manuscripts are in the public
% domain; what this line declares is not a right but a quality — an
% unverified first pass by a machine. The skills write the line; a file
% without it gets no watermark, so leaving it out is visible in the output.
\newcommand{\watermark}[1]{\gdef\grothWatermark{#1}}

\gdef\grothFolder{?}\gdef\grothBatch{}\gdef\grothFirst{?}\gdef\grothLast{?}%
\gdef\grothFoldertitle{}\gdef\grothShelfmark{}\gdef\grothDating{}\gdef\grothWatermark{}

% --- The résumé -------------------------------------------------------------
\newenvironment{resume}
  {\small\setlength{\parskip}{0.45em}}
  {\par\medskip\hrule\medskip}

% --- Keywords ---------------------------------------------------------------
% In English, closing the modernised reading's résumé: the modern vocabulary
% under which to search for what the leaves construct. The manifest extracts
% them to tag the volume — this line is the single source of the tags.
\newcommand{\keywords}[1]{\par\smallskip\noindent
  {\footnotesize\textsc{Keywords} --- \textit{#1}}\par}

% --- The critical apparatus -------------------------------------------------

% The Gallica view number, in the margin. This is the anchor that lets a
% disagreement over a formula be settled by looking at *one* image rather than
% a volume of 750. The site uses it to turn the facsimile as the transcription
% is scrolled. A view is one scanned image: usually one side of a leaf.
\newcommand{\page}[1]{%
  \marginnote{\footnotesize\color{gray!70}\textsf{v.\,#1}}%
  \label{page:#1}\ignorespaces}

% The BnF's pencilled foliation, where the leaf carries it and a pass can read
% it: « 198r ». This is what the literature cites — « ff. 198r–208v » — and
% the one line that turns a citation into a view. Recorded, never inferred:
% a folio a pass cannot read is not written.
\newcommand{\folio}[1]{%
  \marginnote{\footnotesize\color{gray!50}\textsf{f.\,#1}}\ignorespaces}

% The modernised reading's coarser counterpart to \page{}: which views a
% section is drawn from.
\newcommand{\pagerange}[2]{%
  \marginnote{\footnotesize\color{gray!70}\textsf{v.\,#1--#2}}%
  \ignorespaces}

% Illegible: never guessed.
\newcommand{\ill}{\textcolor{red!60!black}{[\,\ldots\,]}}

% A doubtful reading: the word is given, and flagged as uncertain.
\newcommand{\uncertain}[1]{\uline{#1}}

% An editorial addition — a missing letter, an implied word.
\newcommand{\add}[1]{\textcolor{blue!50!black}{[#1]}}

% Struck out by the writer. What was crossed out is part of the document.
\newcommand{\struck}[1]{\sout{#1}}

% The transcriber's note, on the state of the page or the mathematics.
\newcommand{\note}[1]{\par\smallskip{\small\color{gray!60!black}\textit{[#1]}}\par\smallskip}

% A marginal note of hers, distinct from the body of the text.
\newcommand{\marginal}[1]{\par\smallskip\noindent\hspace{1em}%
  {\small\color{gray!60!black}#1}\par\smallskip}

% --- Title ------------------------------------------------------------------
% The title block is French, like the documents themselves.

\AddToHook{shipout/background}{%
  \ifx\grothWatermark\empty\else
    \begin{tikzpicture}[remember picture, overlay]
      \node[rotate=52, scale=3.4, align=center, text=gray!18,
            font=\sffamily\bfseries]
        at (current page.center) {\grothWatermark};
    \end{tikzpicture}%
  \fi}

\AtBeginDocument{%
  \begin{center}
    {\large\bfseries Manuscrits de mathématiciens (Gallica) ---
      \ifx\grothShelfmark\empty\grothFolder\else\grothShelfmark\fi}\\[2pt]
    {\itshape\grothFoldertitle}\\[4pt]
    {\small vues \grothFirst--\grothLast
      \ifx\grothBatch\empty\else\ (lot \grothBatch)\fi}\\[2pt]
    \ifx\grothDating\empty\else
      {\small\color{gray!60!black}Datation du catalogue : \grothDating}\\[2pt]
    \fi
    \ifx\grothWatermark\empty\else
      {\small\bfseries\color{red!45!gray}\grothWatermark}\\[2pt]
    \fi
    {\footnotesize\color{gray!60!black}Transcription établie d'après les images IIIF de Gallica.
     Source gallica.bnf.fr / Bibliothèque nationale de France.\\
     Les lectures incertaines sont soulignées, les illisibles notées [\,\ldots\,],
     les ajouts entre crochets ; \textsf{v.} numérote les vues de Gallica,
     \textsf{f.} la foliotation au crayon de la BnF.}
  \end{center}
  \vspace{1em}}

\endinput


\folder{fr-9115}
\shelfmark{Français 9115}
\batch{11}
\pages{201}{220}
\dating{1801-1900}
\watermark{Lecture automatique\\non vérifiée}
\foldertitle{Papiers de Sophie GERMAIN. — Recueil de dissertations et problèmes mathématiques et physiques.}

\begin{document}

\note{Numérotation des feuillets. Trois séries, toutes à l'encre. La copie
latine porte en tête de chaque page, au milieu, un nombre entre parenthèses,
de (157) à (161) aux vues 201 à 206, qui semble reporter la page de
l'imprimé copié : c'est une inférence. La série à l'encre des lots
précédents se poursuit en haut à droite : 98 (vue 201), 99 (203, et sur sa
seconde photographie, 204), 100 (206), puis, sur les feuillets français,
101 (208), 102 (210), 103 (212), 104 (214), 105 (216), 106 (218), 107 (220),
un par recto écrit et sans lacune. Elle numérote enfin elle-même les
feuillets de chacune de ses deux rédactions : 2 (212), 3 (214), 4 (216) en
haut à gauche pour la première ; 1 (218) et 2 (220), soulignés, en haut au
milieu pour la seconde. Tous ces nombres ont été lus sur la vue ; aucun
n'est déduit. Aucune foliotation au crayon de la Bibliothèque n'a été lue
sur ces vues, et aucun « f. » n'est donc porté ici. La vue 204 photographie
une seconde fois la page de la vue 203 et n'est pas transcrite ; les vues
207, 209, 211, 213, 215, 217 et 219 sont des versos blancs où ne paraît que
la transparence de leur recto.}

\page{201}
(157)

porro
\[ \alpha'=0,\quad \beta'=0,\quad \gamma'=(e^{\omega}-1)(1-e^{\omega}) \text{ et } \delta'=0. \]
Quia igitur omnes evanescunt praeter $\gamma$ et $\gamma'$, ac praeterea est $\gamma'=-\gamma$,
si ponamus $\gamma=1$ erit $\gamma'=-1$; pro portione $EL$ igitur aequatione
motum exprimens erit
\[ y=C\,\text{sin.}\Big(\zeta+t\frac{\omega\omega c\surd 2gb}{aa}\Big)\text{sin.}\,u\omega \]
et pro altera portione $LF$
\[ y=-C\,\text{sin.}\Big(\zeta+t\frac{\omega\omega c\surd 2gb}{aa}\Big)\text{sin.}\,u\omega \]
ubi est $u=\frac{s}{a}$.

§. 51. Praeterea vero datur adhuc alia solutio ex altero factore
oriunda, ex quo fit
\[ \text{tang.}\,\tfrac{1}{2}\omega=\frac{e^{\omega}-1}{e^{\omega}+1}=\frac{e^{\frac{1}{2}\omega}-e^{-\frac{1}{2}\omega}}{e^{\frac{1}{2}\omega}+e^{-\frac{1}{2}\omega}}, \]
qui valor congruit cum eo, qui supra, casu quinto, est erutus:
totum enim discrimen in hoc consistit, quod hic sit $\frac{1}{2}\omega$ quod ubi erat $\omega$,
quemadmodum rei natura postulat, quoniam hic utriusque portionis
longitudo tantum est $\frac{1}{2}a$. Hinc ergo intelligimus, utramque portionem
$EL$ et $LF$ perinde contremiscere posse ac si utraque in $E$ et $F$
stylo simpliciter esset fixa, in $L$ vero firmiter prorsus infixa. Hic
autem valores pro $\omega$ erunt $\frac{5\pi}{2}$, $\frac{9\pi}{2}$, $\frac{13\pi}{2}$, $\frac{17\pi}{2}$, $\frac{21\pi}{2}$, \&c. itaque hinc
soni orientur duabus octavis altiores quam casu quinto. Hic igitur
maxime notatu dignum contigit, quod ambas portiones $EL$ et $LF$
duplici modo contremiscere possunt, altero, qui cum casu superiore
\note{« (157) » en haut, au milieu ; le nombre 98 est à l'encre en haut à droite. Le « et » de la première ligne d'équations est traversé par la barre de son $t$, rejetée en arrière : ce n'est pas une rature. Le radical est écrit sans barre, « $\surd 2gb$ ». Le feuillet s'arrête sur « altero, qui cum casu superiore », dont la phrase se poursuit à la vue 202.}

\page{202}
(158)

quinto, altero vero, qui cum casu quinto congruit. Caeterum valores
coefficientium perinde se habebunt, uti jam supra sunt evoluti,
nisi quod pro hoc casu non fit $\text{sin.}\,\frac{1}{2}\omega=0$

\subsection{Scholion}

§. 52. Hactenus perpetuo assumsimus, virgam elasticam in
statu naturali esse rectam. Nihilo vero difficilior evadit investigatio,
si virga in statu naturali habuerit figuram quamcunque incurvatam.
Veluti si ejus figura naturalis fuerit curva quaecunque $EXF$,
\marginal{Tab III fig 9}
totum discrimen huc reducitur, ut ipsam hanc lineam curvam $EXF$
\add{dum ante nobis erat linea recta cum figura curvae congruens.}
pro axe accipiamus. Hic scilicet sumto portione quacunque
$EX=x=s$, punctum virgae $X$ alium motum recipere nequit,
nisi in directione $XY$, ad ipsam curvam normali. Quare si
concipiamus durante motu punctum $X$ translatum esse $Y$, ac vocemus
\uncertain{hance} applicatam $XY=y$, formulae differentiales, quas theoria
nobis suggessit etiam hic locum habebunt, atque adeo formula
$-\frac{ddy}{ds^2}$, hic jam excessum curvaturae in $Y$ supra curvaturam naturam
in $X$ exprimet: quo observato aequatio motum determinans manebit
prorsus ut ante, sed scilicet
\[ \frac{1}{2g}\left(\frac{ddy}{dt^2}\right)=-bcc\left(\frac{d^4y}{ds^4}\right), \]
unde etiam pro singulis motibus regularibus habebitur eadem
aequatio integralis:
\[ y=C\,\text{sin.}\Big(\zeta+t\frac{\surd 2g}{k}\Big)\Big(\alpha e^{\frac{s}{f}}+\beta e^{-\frac{s}{f}}+\gamma\,\text{sin.}\,\frac{s}{f}+\delta\cos.\frac{s}{f}\Big) \]
existente $f^4=bcck$. Quare si tota virgae longitudo $EXF$ statuatur $=a$,
\note{« (158) » en haut, au milieu ; cette page ne porte pas de nombre en haut à droite. Deux petits crochets à la plume dans la marge gauche, devant « coefficientium » et sous son début. Dans « non fit sin. $\frac{1}{2}\omega=0$ », « fit » est traversé par la barre de son $t$ rejetée en arrière et « sin. » est repassé et empâté ; ce n'est pas une rature. La phrase « dum ante nobis erat linea recta cum figura curvae congruens. » est écrite d'une plume beaucoup plus fine au-dessus de la ligne, précédée d'une croix et rappelée par une croix placée au-dessus de « accipiamus » : elle est mise à cet endroit. « sumto portione » est écrit ainsi, là où la grammaire attendrait « sumta ». Le mot lu « hance » a ses dernières lettres repassées et pourrait n'être que « hanc ». « supra curvaturam naturam » est écrit ainsi. « Scholion » est souligné. Le feuillet s'arrête sur « statuatur $=a$, ».}

\page{203}
(159)

omnes casus, quos supra pertractavimus etiam hic sine ulla mutatione locum
habebunt, et ista virga omnes illos sonos edere poterit, quos supra assignavimus,
prouti scilicet virgae termini fuerint liberi, vel simpliciter fixi, vel etiam
firmiter infixi; ita ut pro talibus virgis naturaliter incurvatis nulla
nova investigatione sit opus. Interim tamen hinc ii casus sunt
excipiendi, quibus virga ita est incurvata, ut figuram in se redeuntem
referat, quando quidem similis casus in virgis naturaliter rectis
locum habere nequit, quamobrem istum casum coronidis loco
hic subjungamus.

\section{Evolutio casus quo virga elastica in statu naturali figuram in se redeuntem habet, sive de sonis annulorum elasticorum}

§. 53. Sit igitur figura virgae elasticae circulus $AXBC$, sive alia
\marginal{Tab. III fig 10}
quaecunque curva in se rediens, cujus tota peripheria sit $=a$,
crassities vero et elasticitas virgae maneant eadem ut ante stabilitae;
tum vero pro motibus regularibus, quos haec virga recipere potest,
sit \emph{k} longitudo penduli simplicis isochroni, unde formetur quantitas
$f=\sqrt[4]{bcck}$; tum pro quacunque portione indefinita $AX=s$ sit $Y$ punctum,
in quod praesenti tempore $=t$ punctum $X$ sit translatum, et jam
vidimus, aequationem integralem in genere fore,
\[ y=C\,\text{sin.}\Big(\zeta+t\frac{\surd 2g}{k}\Big)\Big(\alpha e^{\frac{s}{f}}+\beta e^{-\frac{s}{f}}+\gamma\,\text{sin.}\,\frac{s}{f}+\delta\cos.\frac{s}{f}\Big). \]
\note{« (159) » en haut, au milieu ; le nombre 99 est à l'encre en haut à droite. Entre « virga » et « omnes », un trait vertical à la plume, comme aux vues 182 et suivantes. « quamobrem » a son $o$ repassé et empâté ; « coronidis » est empâté de même. Le titre est écrit sur trois lignes, « Evolutio casus » souligné et la dernière ligne soulignée d'un long trait. Le radical de $f$ est écrit sans barre, surmonté d'un petit chiffre lu « 4 » : il est rendu $\sqrt[4]{bcck}$, et le $k$ final est une capitale barrée. Les deux chiffres qui suivent « fig » en marge sont empâtés ; la lecture « 10 » suit la suite des figures de la planche, la vue 202 citant « fig 9 ». Le $A$ de « $AX=s$ » est surchargé d'une tache. Le feuillet s'arrête sur cette équation ; la suite est à la vue 205. Cette page est photographiée une seconde fois à la vue 204.}

\page{205}
(160)

Sicque totum negotium jam huc perductum, ut coefficientibus
$\alpha$, $\beta$, $\gamma$, $\delta$ debiti valores assignentur, ubi res longe aliter se habere
deprehenditur ac supra, quoniam hic neque termini liberi, neque
simpliciter fixi, neque firmiter infixi occurrunt.

§. 54. Est vero ipsa indoles, qua figura virgae in se
rediens assumitur facilem viam nobis aperit hos coefficientes
determinandi. Consideremus enim ipsum punctum $A$, ubi est
$s=0$, quod ita accipi potest, ut ibi fiat etiam $y=0$. Pro
hoc ergo puncto, omisso factore partim constante partim a
tempore pendente, fieri debet $0=\alpha+\beta+\delta$. Jam statuamus
$s=a$, et quia iterum in idem punctum $A$ incidimus ponendo
brevitatis gratia $\frac{a}{f}=\omega$, fieri oportet
\[ 0=\alpha e^{\omega}+\beta e^{-\omega}+\gamma\,\text{sin.}\,\omega+\delta\cos\omega, \]
atque idem evenire debet, si ponamus $s=2a$ sive $3a$
sive $4a$ \&c unde nascentur sequentes aequationes:
\[ 0=\alpha e^{2\omega}+\beta e^{-2\omega}+\gamma\,\text{sin.}\,2\omega+\delta\cos.2\omega, \]
\[ 0=\alpha e^{3\omega}+\beta e^{-3\omega}+\gamma\,\text{sin.}\,3\omega+\delta\cos 3\omega, \]
\[ \&c \qquad \&c. \]
quibus omnibus simul satisfieri nequit, nisi sit $\alpha=0$, $\beta=0$, et
$\delta=0$; tum vero necesse est, ut simul fiat
\[ \text{sin.}\,\omega=0,\quad \text{sin.}\,2\omega=0,\quad \text{sin.}\,3\omega=0 \ \&c. \]
quod in genere evenit, sumendo $\omega=i\pi$, denotante $i$ numerum
\note{« (160) » en haut, au milieu ; le « 6 » est repassé et empâté. Cette page ne porte pas de nombre en haut à droite. Entre « res » et « longe », un trait vertical à la plume. Les deux « \&c » sous les deux dernières équations sont écrits l'un sous $\beta e^{-3\omega}$, l'autre sous $\delta\cos 3\omega$ ; ils sont rendus sur une ligne. Le feuillet s'arrête sur « denotante $i$ numerum ».}

\page{206}
(161)

integrum quemcunque; unde ut supra oritur sonus $=\frac{ii\pi c\surd 2gb}{aa}$, ita ut soni
simplices progrediantur in hac progressione.
\[ \frac{\pi c\surd 2gb}{aa},\ \frac{4\pi c\surd 2gb}{aa},\ \frac{9\pi c\surd 2gb}{aa},\ \frac{16\pi c\surd 2gb}{aa}\ \&c. \]

§. 55. Sonus igitur principalis, quem talis virga edere potest, continetur
in formula $\frac{\pi c\surd 2gb}{aa}$; aliqui vero \struck{progrediuntur, qui cum} soni simplices
secundum numeros quadratos $4, 9, 16, 25$, progrediuntur, qui cum mox
nimis fiant alti quam ut exaudiri queant, praeter sonum fundamentalem
plerumque alius non sentietur, nisi dupla octava, quo ipso harmonia
gratissima \struck{\ill{}} \add{percipietur}. Tales igitur annuli elastici prae chordis musicis
hac insigni proprietate sunt praediti, ut sonos multo puriores reddant,
id quod etiam in integris discis et catinis campaniformibus evenire
debere videtur, cujusmodi corpora inter instrumenta musica jam sunt
recepta, quorumque soni singulari suavitate sensum auditus
afficere feruntur. Caeterum manente eadem elasticitate hi soni
tenent rationem reciprocam duplicatam totius perimetri $a$, prorsus
uti jam supra circa virgas rectas observavimus

fin.
\note{« (161) » en haut, au milieu ; le nombre 100 est à l'encre en haut à droite. « progrediuntur, qui cum » est biffé d'un long trait. Le mot biffé avant « percipietur » est couvert d'une rature épaisse : on y devine « pro… » au début et « …etur » à la fin, mais il n'est pas lu ; « percipietur » est écrit au-dessus, sans signe d'insertion, et mis à sa place. La copie s'arrête ici : après « observavimus », le mot français « fin. » est écrit au milieu de la page, suivi d'un grand paraphe. Le bas du feuillet est blanc.}

\page{208}

\section{Remarques sur la lame elastique}
\note{le feuillet ne porte que ce titre, écrit en grand au haut de la page, et le nombre 101 à l'encre en haut à droite ; tout le reste est blanc. « elastique » est écrit sans accent sur l'$e$ initial, le point de l'$i$ étant reporté sur les lettres qui précèdent. C'est un feuillet de titre pour ce qui suit ; son verso, vue 209, est blanc.}

\page{210}

\section{Remarque sur le mémoire d'Euler: Investigatio motuum quibus lamina et virga elastica contremiscunt, in. act. Acad. pétrop ann. 1779 P.I p. 103 et seq.}

Après avoir examiné les mouvemens possibles dela lame vibrante en ayant seulement
egard a l'état des extremités qui peuvent etre fixées, simplement appuiées ou libres,
d'ou il resulte six cas essentiellement differens; l'auteur suppose que la même lame
puisse en outre etre appuiée dans un ou plusieurs points differens des extremités
demanière a ceque la communication entre ses differentes parties ne soit pas
interrompue. Puis il ajoute que le developpement detous ces cas demanderoit des
recherches infinies \add{et} qu'il suffit d'examiner un d'entr'eux pour faire comprendre
comment le calcul peut leur etre appliqué.

Le cas qu'il choisit est le plus simple de tous \struck{savoir celui} savoir celui ou
la lame elastique est non seulement, simplement appuiée a l'une et l'autre de ses
extremités mais encore soumise a l'action d'un stilet dans un point quelconque
different des extremités il se propose (N\textsuperscript{o} 47) de trouver tous les mouvemens
de vibration dont une telle lame peut etre agitée.

Dans la suite du même N\textsuperscript{o} notre auteur parvient a l'équation suivante:
\[ 0=2-2e^{2\omega}-(e^{\lambda\omega}-e^{\omega(2-\lambda)})(e^{\lambda\omega}-e^{-\lambda\omega})\frac{\text{sin.}\,\omega}{\text{sin}\,\lambda\omega\,\text{sin}(1-\lambda)\omega} \]
il remarque que la question est réduite a tirer de cette équation les valeurs de
l'angle $\omega$. mais il ajoute qu'il n'ose entreprendre ce travail acause de la
complication des formules et qu'il suffit d'avoir montré dans ce lieu par quelle
methode on peut traiter ces difficiles questions.

Cependant après avoir etabli que les suppositions $\lambda=0$ et $\lambda=1$ menent
a conclure deson équation finale cidessus, $0=2(1-e^{2\omega})$ d'ou il suit $1=e^{2\omega}$
et parconséquent $\omega=0$ qui est la valeur qui convient aurepos; il ajoute, que
le cas ou $\lambda=\frac{1}{2}$ mérite d'être developpée.

C'est cequ'il effectue N\textsuperscript{o} 48 et il trouve que son équation finale prend
alors la forme suivante:
\[ 0=2-2e^{2\omega}-(e^{\frac{1}{2}\omega}-e^{\frac{3}{2}\omega})(e^{\frac{1}{2}\omega}-e^{-\frac{1}{2}\omega})\frac{\text{sin}\,\omega}{\text{sin.}^2\frac{1}{2}\omega} \]
qui peut être réduite a celle ci: $0=2(1-e^{2\omega})\,\text{sin}^2\frac{1}{2}\omega-(1-e^{\omega})(e^{\omega}-1)\,\text{sin}\,\omega$,
laquelle devient après ladivision par le facteur commun $1-e^{\omega}$,
$0=2(1+e^{\omega})\,\text{sin}^2\frac{1}{2}\omega-(e^{\omega}-1)\,\text{sin.}\,\omega$ il met ensuite dans cette dernière
équation pour $\text{sin}\,\omega$ sa valeur $2\,\text{sin}\,\frac{1}{2}\omega\cos\frac{1}{2}\omega$ et elle devient après avoir
été divisée par 2, $0=(1+e^{\omega})\,\text{sin}^2\frac{1}{2}\omega-(e^{\omega}-1)\,\text{sin}\,\frac{1}{2}\omega\cos\frac{1}{2}\omega$
laquelle a manifestement deux facteurs l'un $\text{sin}\,\frac{1}{2}\omega$ et l'autre
$(1+e^{\omega})\,\text{sin}\,\frac{1}{2}\omega-(e^{\omega}-1)\cos\frac{1}{2}\omega$ qui l'un et l'autre egalés a zéro donne
une solution. L'auteur examine chacune de ces solutions enparticulier

Il prend d'abord (N. 49) $\text{sin}\,\frac{1}{2}\omega=0$ il en conclut $\frac{1}{2}\omega=i\pi$ en designant par
$i$ un nombre entier quelconque et il observe que ce cas est parfaitement semblable
au cas quatrieme qui est celui ou les deux extremités dela lame vibrante sont
simplement appuiées (\uncertain{2d} N\textsuperscript{o} 35), si ce n'est que l'on a ici $\frac{1}{2}\omega$ a la place
de $\omega$ dans l'endroit cité, et cela parceque l'une et l'autre parties de la longueur
sont seulement $\frac{1}{2}a$.

Il conclut que l'une et l'autre « moitiés » font leurs vibrations delamême manière
que si elles avoient
\note{le nombre 102 est à l'encre en haut à droite ; cette page ne porte pas de nombre entre parenthèses, la copie latine s'étant arrêtée à la vue 206. En haut à gauche, un signe à l'encre, noyé dans une tache triangulaire et suivi du même trait oblique qui accompagne les chiffres 2, 3 et 4 des vues 212, 214 et 216 : c'est la place du premier nombre de la série, mais le chiffre lui-même n'est pas lu. Le titre est souligné d'un bout à l'autre, les mots latins soulignés en outre chacun séparément, et un double filet le sépare du texte. « et » est ajouté d'une plume plus fine au-dessus de la ligne, entre « infinies » et « qu'il ». Un petit signe à la plume dans la marge gauche, en regard de la seconde équation. Dans la référence lue « (2d N\textsuperscript{o} 35) », les deux premiers signes sont un chiffre 2 suivi d'une lettre bouclée qui n'est pas assurée. Le mot « moitiés » est entouré sur le feuillet de deux petites virgules hautes, rendues ici par des guillemets. L'orthographe est la sienne : « mouvemens », « differens », « appuiées », « demanière », « detous », « demanderoit », « acause », « parconséquent », « aurepos », « cidessus », « parceque », « delamême », « dela », « deson », « enparticulier », « ladivision », « cequ'il ». Le feuillet s'arrête sur « que si elles avoient » ; la suite est à la vue 212.}

\page{212}

leurs extremités appuiées au point ou est le stilet d'ou il resulte que les sons simples
que l'une et l'autre portions peuvent donner sont exprimés par les nombres suivans:
\[ \frac{4\pi c\surd 2gb}{aa};\ \frac{16\pi c\surd 2gb}{aa};\ \frac{36\pi c\surd 2gb}{aa};\ \frac{64\pi c\surd 2gb}{aa}; \]
et sont parconséquent les doubles octaves audessus, de ceux qui ont lieu dans le cas IV.
cequi est naturel puisque dans le cas present les longueurs sont seulement la moitié
de celles du cas IV.

L'auteur etablit ensuite N\textsuperscript{o} 50 que les coëfficiens $\alpha$, $\beta$, $\gamma$, $\delta$ et
$\alpha'$, $\beta'$, $\gamma'$, $\delta'$ seront déterminés pour le cas present dela manière suivante:
\[ \alpha=e^{\frac{1}{2}\omega}-e^{\frac{3}{2}\omega};\quad \beta=e^{\frac{3}{2}\omega}-e^{\frac{1}{2}\omega},\quad \gamma=-\frac{(e^{\frac{1}{2}\omega}-e^{\frac{3}{2}\omega})(e^{\frac{1}{2}\omega}-e^{-\frac{1}{2}\omega})}{\text{sin.}\,\frac{1}{2}\omega};\quad \delta=0 \]
\[ \alpha'=e^{\frac{1}{2}\omega}-e^{-\frac{1}{2}\omega};\quad \beta'=-e^{2\omega}(e^{\frac{1}{2}\omega}-e^{-\frac{1}{2}\omega});\quad \gamma'=-\frac{(e^{\frac{1}{2}\omega}-e^{-\frac{1}{2}\omega})(e^{\frac{1}{2}\omega}-e^{\frac{3}{2}\omega})}{\text{sin.}\,\frac{1}{2}\omega-\cos\frac{1}{2}\omega\,\text{tang}\,\omega} \]
\[ \text{et}\quad \delta'=\frac{(e^{\frac{1}{2}\omega}-e^{-\frac{1}{2}\omega})(e^{\frac{1}{2}\omega}-e^{\frac{3}{2}\omega})\,\text{tang}\,\omega}{\text{sin}\,\frac{1}{2}\omega-\cos\frac{1}{2}\omega\,\text{tang}\,\omega} \]
et il conclut acause de
$\text{sin}\,\frac{1}{2}\omega-\cos\frac{1}{2}\omega\,\text{tang}\,\omega=\frac{\text{sin}\frac{1}{2}\omega\cos\omega-\cos\frac{1}{2}\omega\,\text{sin}\,\omega}{\cos\omega}=-\frac{\text{sin}\,\frac{1}{2}\omega}{\cos.\omega}$
et de $\omega=2i\pi$ d'ou il resulte $\cos\omega=1$, et $\text{sin}\,\frac{1}{2}\omega=0$; que si on multiplie tous
ces coëfficiens par $\text{sin}\,\frac{1}{2}\omega$ ils deviendront
\[ \alpha=0,\ \beta=0,\ \gamma=-(1-e^{\omega})(e^{\omega}-1),\ \delta=0 \text{ et } \alpha'=0,\ \beta'=0,\ \gamma'=(e^{\omega}-1)(1-e^{\omega}),\ \delta'=0. \]

On voit en effet qu'il est permis de multiplier tous ces coëfficiens par $\text{sin}\,\frac{1}{2}\omega$
ou par tel autre facteur que l'on voudra puisque les valeurs d'$\alpha$ et $\alpha'$ desquelles
resultent celles des autres coëfficiens sont conclues dela fraction
$\frac{\alpha}{\alpha'}=\frac{e^{\lambda\omega}-e^{(2-\lambda)\omega}}{e^{\lambda\omega}-e^{-\lambda\omega}}$. (Voyez N\textsuperscript{o} 47)

Le N\textsuperscript{o} 51 est consacré a l'examen de la solution qui resulte de l'autre facteur
duquel on tire $\text{tang}\,\frac{1}{2}\omega=\frac{e^{\omega}-1}{1+e^{\omega}}=\frac{e^{\frac{1}{2}\omega}-e^{-\frac{1}{2}\omega}}{e^{\frac{1}{2}\omega}+e^{-\frac{1}{2}\omega}}$ l'auteur observe que cette
valeur est lamême que celle qu'il a trouvé pour le cas V ou une extremité est appuiée
et l'autre fixée et que la seule difference consiste enceque l'on a ici $\frac{1}{2}\omega$ au lieu
de $\omega$, cequi est naturel puisque l'une et l'autre portions dela longueur sont
seulement $\frac{1}{2}a$. Il conçoit que les deux moitiés dela lame puissent se mouvoir
en même tems et qu'ayant chacune une deleurs extrèmités appuiées la seconde
se trouve fixée par le stilet consideré non plus comme simple point d'appui mais
comme un obstacle fixe. Il trouve que les valeurs de l'angle $\omega$ doivent être
$\frac{5\pi}{2}$, $\frac{9\pi}{2}$, $\frac{13\pi}{2}$ \&c. et que le ton est parconséquent de deux octaves plus élevé
que dans le cas V.

Il ajoute enfin qu'il est digne de remarque que les deux moitiés dela lame
puissent vibrer de deux manières, l'une qui se rapporte au cas IV et l'autre
qui appartient au cas V.

Il me semble étonnant qu'Euler ait pû croire que cette seconde solution satisfait
a laquestion puisque son hypothèse est générale, dont alavérité il n'examine
qu'un seul cas et: qu'independamment de l'état des extrèmités la lame puisse
\uncertain{pu} etre appuiée dans un ou plusieurs points de manière a ceque la
communication entre ses differens points ne soit pas interrompue. et que la
supposition d'un point fixe aumilieu de l'étendue dela même lame équivaut
a sa séparation en deux moitiés entièrement independantes.

D'ailleurs on verra bientôt que pour toute autre valeur de $\lambda$ que $\frac{1}{2}$ il ne se
trouve point defacteur de l'équation finale d'ou on puisse conclure une solution dela
\note{le nombre 103 est à l'encre en haut à droite ; en haut à gauche, un chiffre 2, de sa main, qui paraît numéroter les feuillets de la Remarque. « enfin », dans « Il ajoute enfin », est repassé sur une première forme empâtée. Le passage « qu'independamment de l'état des extrèmités … ne soit pas interrompue » est souligné d'un bout à l'autre. Entre « la lame puisse » et « etre appuiée », un mot court de deux ou trois lettres, lu « pu », qui ne fait pas sens ; il n'est pas biffé. Les longs traits horizontaux qui traversent « extremité est », « toute autre » et d'autres mots sont les barres des $t$ rejetées en arrière, non des ratures. Dans $\beta'$, elle écrit $e^{\frac{1}{2}\omega}-e^{-\frac{1}{2}\omega}$ là où la vue 200 écrivait deux fois le même exposant. Le feuillet s'arrête sur « une solution dela » ; la suite est à la vue 214.}

\page{214}

question relative au cas V. tandis qu'au contraire la solution qui serapporte au cas IV
est générale et independante de la valeur de $\lambda$ pourvu toutefois que $\lambda$ soit une
partie aliquote de l'unité.

Pour se rendre raison dela difficulté qui semble résulter de l'impossibilité d'admettre
comme solution dela question la \struck{condition} \add{valeur}
$\text{tang}\,\frac{1}{2}\omega=\frac{e^{\frac{1}{2}\omega}-e^{-\frac{1}{2}\omega}}{e^{\frac{1}{2}\omega}+e^{-\frac{1}{2}\omega}}$ qui dans le cas particulier
ou $\lambda=\frac{1}{2}$ satisfait cependant a l'équation finale, je crois qu'il suffit de remarquer
que la condition dela communication entre les differens points dela lame ne passe
pas dans le calcul, d'ou il resulte que l'équation finale peut avoir des solutions
qui soient independantes de cette condition et qu'il faut rejetter si on veut y
avoir egard pour ne conserver quecelles qui sont compatibles avec elle. Il semble
\struck{en effet} d'ailleurs que l'examen de l'effet des points d'appuis appliqués dans
des lieux differens des extremités ne doive jamais donner d'autres resultats que
ceux qui sont déja contenus dans la solution générale qui renferme tous les
mouvemens possibles dela lame pour un etat donné des extremités, car cet effet
est consideré comme purement exclutif et puisqu'il est évident qu'il ne peut faire
naître des mouvemens inconciliables avec l'état des extrèmités il se borne a
déterminer quelques uns des mouvemens possibles déja connus en rendant les
autres impossibles; cela posé il faut que toute solution dela question soit
conciliable avec la solution générale relative a l'état des extremités et on voit
que dans le cas present qui est celui des deux extremités appuiées l'angle $\omega$
est déterminé par l'équation $\text{sin.}\,\omega=0$. cette équation est d'accord avec celle
que fournit l'équation finale pour le cas qui se rapporte au cas V savoir
$\text{sin}\,\frac{1}{2}\omega=0$ d'ou resulte
\marginal{$2\,\text{sin}\,\frac{1}{2}\omega\cos\frac{1}{2}\omega=\text{sin}\,\omega=0$}
mais il n'en est pas demême dela valeur $\text{tang}\,\frac{1}{2}\omega=\frac{e^{\omega}-1}{e^{\omega}+1}$ d'ou on tire
\marginal{$\frac{1}{\cos^2\frac{1}{2}\omega}=\frac{2(e^{2\omega}+1)}{(e^{\omega}+1)^2}$}
\marginal{$\cos^2\frac{1}{2}\omega=\frac{(e^{\omega}+1)^2}{2(e^{2\omega}+1)}$}
\marginal{$1-\cos^2\frac{1}{2}\omega=\frac{(e^{\omega}-1)^2}{2(e^{2\omega}+1)}$}
\marginal{$\text{sin}\,\omega=\uncertain{2}\frac{(e^{2\omega}-1)}{2(e^{2\omega}+1)}$}
car on est mené a conclure $0=\frac{e^{2\omega}-1}{2(e^{2\omega}+1)}$ ou $e^{2\omega}=1$, $\omega=0$. c'est adire
que l'on tombe sur lavaleur qui rend le mouvement impossible puisqu'elle
convient au repos absolu delalame.

Au reste on verra dans la suite que si on seproposé d'examiner l'effet de l'action
d'un stilet appliqué a un point different des extremités en supposant les mêmes
extrèmités l'une et l'autre libres et ensuite l'une et l'autre fixées on parvient
toujours a une équation finale qui dans le cas ou $\lambda=\frac{1}{2}$ donne deux solutions
l'une relative au cas ou la lame seroit consideré comme separée en deux
lames moitiés moins longues dont une des extremités seroit libre, si les deux
\note{le nombre 104 est à l'encre en haut à droite ; en haut à gauche, le chiffre 3, de sa main. « valeur » est écrit au-dessus de « condition » biffé. « en effet » est biffé d'un trait épais, « d'ailleurs » le suit sur la ligne. « naître » est repassé sur une première forme et traversé par la barre de son $t$ : il n'est pas biffé. « exclutif » est écrit ainsi. Les cinq formules données ici comme notes marginales sont écrites dans la marge de droite du feuillet, en regard des lignes où elles sont placées ; la dernière porte, entre le signe d'égalité et la fraction, un signe empâté lu « 2 » — avec ce facteur le calcul se referme, et c'est la seule raison de le lire ainsi. Le feuillet s'arrête sur « si les deux » ; la suite est à la vue 216.}

\page{216}

extremités dela lame primitive ont été supposées libres, fixe si les mêmes extremités
ont été supposées fixes, tandis que la seconde extremité seroit appuiée aupoint du
milieu ou le stilet est appliqué, et l'autre solution relative au cas ou tout restant
demême que dans \struck{l'autre} la premiere solution l'extremité que l'on vient de
supposer appuiée au point ou le stilet est appliqué seroit fixée au même point.
Et quelemême que dans le cas calculé par Euler, qui est celui ou les extremités dela
lame primitive sont l'une et l'autre appuiées, la premiere solution est la seule
qui s'accorde avec la valeur generale de l'angle $\omega$, tandis que la seconde
mene a conclure $\omega=0$ et n'appartient pas parconséquent a laquestion dans
laquelle on seproposé de déterminer les mouvemens possibles de la lame soumise
aux conditions donnees, puisqu'elle exige que lamême lame reste en repos.

Enfin si il etoit necessaire d'ajouter une nouvelle preuve a l'impossibilité de
considerer la seconde solution d'Euler comme appartenant a laquestion il suffiroit
d'observer qu'excepté le cas $\lambda=\frac{1}{2}$, les deux portions dela lame separées
par un appui fixe ne peuvent rendre lemême ton, comme il est aisé de s'en assurer
par la consideration des differens nombres qui leur conviennent, tandis que toute
cette analyse est fondée sur la supposition que les deux portions enquestion se
meuvent dans lemême tems et donnent parconséquent des sons pareils

Je reprens apresent l'équation finale
$0=2-2e^{2\omega}-(e^{\lambda\omega}-e^{\omega(2-\lambda)})(e^{\lambda\omega}-e^{-\lambda\omega})\frac{\text{sin}\,\omega}{\text{sin}\,\lambda\omega\,\text{sin}(1-\lambda)\omega}$
donné par Euler \struck{de} de laquelle il s'agit de tirer les valeurs de l'angle $\omega$,
et je me propose de prouver que quelque soit la valeur de $\lambda$ pourvu toutefois
que cette quantité soit fraction exacte de l'unité, l'équation finale est toujours
susceptible d'une solution qui se rapporte au cas quatrieme.

Cette équation
$0=2-2e^{2\omega}-(e^{\lambda\omega}-e^{(2-\lambda)\omega})(e^{\lambda\omega}-e^{-\lambda\omega})\frac{\text{sin}\,\omega}{\text{sin}\,\lambda\omega\,\text{sin}(1-\lambda)\omega}$
peut être transformée comme il suit:
\[ 0=2(e^{-\omega}-e^{\omega})-(e^{(\lambda-1)\omega}-e^{(1-\lambda)\omega})(e^{\lambda\omega}-e^{-\lambda\omega})\frac{\text{sin}\,\omega}{\text{sin}\,\lambda\omega\,\text{sin}(1-\lambda)\omega} \]
\[ 0=\frac{2(e^{\omega}-e^{-\omega})}{(e^{(\lambda-1)\omega}-e^{(1-\lambda)\omega})(e^{\lambda\omega}-e^{-\lambda\omega})}\cdot \text{sin}\,\lambda\omega\,\text{sin}(1-\lambda)\omega \ +\ \text{sin}\,\omega \]
\begin{align*}
0={}& \frac{(e^{(\lambda-1)\omega}+e^{(1-\lambda)\omega})(e^{\lambda\omega}-e^{-\lambda\omega})+(e^{\lambda\omega}+e^{-\lambda\omega})(e^{(1-\lambda)\omega}-e^{(\lambda-1)\omega})}{(e^{(\lambda-1)\omega}-e^{(1-\lambda)\omega})(e^{\lambda\omega}-e^{-\lambda\omega})}\,\text{sin}\,\lambda\omega\,\text{sin}(1-\lambda)\omega \\
&+ \text{sin}\,\lambda\omega\cos(1-\lambda)\omega + \text{sin}(1-\lambda)\omega\cos\lambda\omega
\end{align*}
\begin{align*}
0={}& \frac{e^{(\lambda-1)\omega}+e^{(1-\lambda)\omega}}{e^{(\lambda-1)\omega}-e^{(1-\lambda)\omega}}\cdot\text{sin}\,\lambda\omega\,\text{sin}(1-\lambda)\omega - \frac{e^{\lambda\omega}+e^{-\lambda\omega}}{e^{\lambda\omega}-e^{-\lambda\omega}}\,\text{sin}\,\lambda\omega\,\text{sin}(1-\lambda)\omega \\
&+ \text{sin}(1-\lambda)\omega\cos\lambda\omega + \text{sin}\,\lambda\omega\cos(1-\lambda)\omega
\end{align*}
enfin
\begin{align*}
0={}& \text{sin}\,\lambda\omega\Big\{\frac{e^{(1-\lambda)\omega}+e^{(\lambda-1)\omega}}{e^{(1-\lambda)\omega}-e^{(\lambda-1)\omega}}\,\text{sin}(1-\lambda)\omega-\cos(1-\lambda)\omega\Big\} \\
&+ \text{sin}(1-\lambda)\omega\Big\{\frac{e^{\lambda\omega}+e^{-\lambda\omega}}{e^{\lambda\omega}-e^{-\lambda\omega}}\,\text{sin}\,\lambda\omega-\cos\lambda\omega\Big\}
\end{align*}
sous cette forme on voit aisément que l'équation finale est fonction homogène
de $\lambda$ et $1-\lambda$
\note{le nombre 105 est à l'encre en haut à droite ; en haut à gauche, le chiffre 4, de sa main. Le mot biffé avant « la premiere solution » est « l'autre » ; celui biffé après « Euler » est « de », aussitôt repris. « Et quelemême » est écrit d'un seul trait, là où le sens demande « Et que de même » ; laissé tel quel. « seproposé », « donnees », « lamême », « lemême », « enquestion », « laquestion », « apresent » sont écrits ainsi. Le mot « enfin » est placé dans la marge gauche, en regard de la dernière équation. Le feuillet s'arrête sur « de $\lambda$ et $1-\lambda$ » ; le bas de la page est blanc et la suite est à la vue 218.}

\page{218}

\section{Remarques sur le memoire d'Euler: Investigatio motuum quibus laminæ et virgæ elasticæ contremiscunt, in act. Acad. Petrop. pro anno 1779 P.I p. 103 seq.}

\struck{\ill{}}. Euler avoit déja traité les questions qui font l'objet de ce memoire
dans sa methodus inveniendi curvas maximi minimique gaudentes,
imprimée en 1744. En reprenant cette matière dans le mémoire dont
il s'agit ici il a exposé dans un ordre plus méthodique les recherches
qui s'appliquent aux six cas qu'il distingue par rapport a l'état
des extremités dela lame elastique vibrante. mais on ne remarque
aucun changement essentiel dans sa méthode et les resultats sont
les mêmes de part et d'autre. L'addition la plus importante
qu'il ait faite a son premier travail concerne l'examen de l'effet
d'un point d'appui placé dans un lieu quelconque de l'étendue
dela lame different des extremités. Je vais rapporter ici cette
partie de son mémoire en y joignant les problèmes relatifs aux
cas qu'il a négligé de traiter.
\marginal{et les reflexions que \struck{l'étude} l'examen plus aprofondi de cette matière m'a fait naitre}
Voici comment s'exprime l'auteur \uncertain{p. 151}

« nous n'avons encore examiné que les cas ou les extremités sont libres,
« appuiées, ou fixés, mais on peut supposer que la lame est appuiée « dans
un ou plusieurs points differens des extremités de manière que « la
communication du mouvement entre ses differentes parties ne soit « pas
interrompue. mais le developpement de tous ces cas demanderoit « des
recherches infinies. C'est pourquoi il suffit d'examiner un seul « pour faire
comprendre comment le calcul peut leur etre « appliqué

« \emph{Probleme}

« 47 Si la lame elastique est non seulement simplement appuiée a l'une « et
l'autre de ses extremités $E$. et $F$., mais encore soumise a l'action « d'un
stilet dans un point quelconque $L$, de son etendue ; trouver « tous les
mouvemens de vibration dont cette lame peut être agitée

« \emph{Solution}

« En continuant de prendre la longueur totale de la lame « $EF=a$ et faisant
la portion $EL=\lambda a$, desorte que $\lambda$ puisse « representer une
fraction quelconque moindre que l'unité. Il est « évident d'abord que si la
lame est fermement fixée au point « $L$. la communication entre les deux
portions est
\note{le nombre 106 est à l'encre en haut à droite ; en haut, au milieu, un petit trait souligné, apparemment le chiffre 1 : ce feuillet ouvre une seconde rédaction des Remarques, dont la vue 210 donne la première. Le titre est souligné, ses mots latins soulignés en outre chacun, et un double filet le sépare du texte. Le premier mot du texte, une abréviation de deux lettres suivie d'un point, est couvert d'une rature en tire-bouchon et n'est pas lu. L'addition marginale, dans la marge de gauche, est précédée d'une croix et porte une seconde croix en exposant après « que » ; aucun signe d'insertion n'a été trouvé dans la ligne, et elle est placée ici d'après le sens. Les chiffres qui suivent « l'auteur », en petit, se lisent « 151 » ou « 131 », le chiffre du milieu n'étant pas assuré. À partir de « nous n'avons », chaque ligne de la citation s'ouvre d'un guillemet répété, à la manière du temps, comme aux vues 138 et 142 ; ces guillemets sont conservés à leur place dans le texte suivi. « Probleme » et « Solution » sont soulignés. Le feuillet s'arrête sur « entre les deux portions est » ; la suite est à la vue 220.}

\page{220}

« entierement interrompue, desorte que le mouvement imprimé « a l'une des
deux ne pourra jamais être communiqué a l'autre « Mais si le stilet est
simplement appuié au point $L$, desorte « que la lame puisse tourner autour
de ce point, ni l'un ni l'autre « de ses portions ne peut recevoir de
mouvement, sans qu'il se « communique d'une manière quelconque a la lame
entière. « Cependant l'action de ce stilet interromp la continuité dela
« courbe entre les deux portions, desorte que la portion $EmL$ « est exprimée
par une autre equation que la portion $LnF$; « desorte que suivant les
principes que nous avons trouvé pour « les mouvemens réguliers, qui sont
conformes au pendul simple « $=k$, le premier facteur
$c\,\text{sin}\big(\zeta+t\surd\frac{g}{k}\big)$ doit necessairement « demeurer
lemême pour l'une et l'autre portions, parceque « les vibrations de ces
portions doivent s'exécuter et se terminer « dans le même tems et delamême
manière. Les autres « facteurs peuvent être differens en raison des
coëfficiens « $\alpha$, $\beta$, $\gamma$, $\delta$. En omettant donc le
premier facteur et fesant, « comme jusqu'a présent $\frac{a}{f}=\omega$ et
$\frac{s}{a}=u$, nous aurons pour « la portion $EL$
\[ y=\ldots\ \alpha e^{u\omega}+\beta e^{-u\omega}+\gamma\,\text{sin}\,u\omega+\delta\cos u\omega \]
« et pour la portion $LF$,
\[ y=\ldots\ \alpha' e^{u\omega}+\beta' e^{-u\omega}+\gamma'\,\text{sin}\,u\omega+\delta'\cos u\omega \]
« dont les coëfficiens peuvent differer des premiers d'une manière
« quelconque, on doit observer seulement, que pour le point $L$, dans lequel
« $u=\lambda$ l'une et l'autre formules doivent donner les mêmes valeurs tant
« pour $\left(\frac{dy}{ds}\right)$ que pour $\left(\frac{ddy}{ds^2}\right)$,
parceque les angles que l'une et « l'autre portions font avec l'axe en $L$,
doivent necessairement etre egaux ; et aussi dans ce point ce point $L$ les
rayons « osculateurs de l'une et l'autre portions ne peuvent etre
« differens. Cela posé, puisque dans ce point $L$, pour lequel « $u=\lambda$,
l'une et l'autre appliquées y doivent s'évanouir, « nous aurons les quatre
équations suivantes.
\note{le nombre 107 est à l'encre en haut à droite ; en haut, au milieu, le chiffre 2 souligné, qui poursuit la numérotation ouverte à la vue 218. Chaque ligne de la citation s'ouvre d'un guillemet répété, conservé ici à sa place dans le texte suivi ; la ligne qui commence par « etre egaux » est la seule qui n'en porte pas. Les deux formules de $y$, écrites en ligne sur le feuillet à la suite de « la portion $EL$ » et de « la portion $LF$ », sont mises ici hors ligne. Dans la formule du premier facteur, le radical est écrit sans barre au-dessus de $\frac{g}{k}$. Les points de conduite après « $y=$ » marquent le facteur omis, comme aux vues 186, 189 et 193. Devant l'$\alpha'$ de la seconde formule, deux ou trois signes sont repassés et empâtés. « dans ce point ce point $L$ » est écrit ainsi, sans rature. Le feuillet s'arrête sur « les quatre équations suivantes. » ; la suite est à la vue 221, hors de ce lot.}

\end{document}
